\chapter{Poles, Zeros, Lasers, and Absorbers}
\label{ch:poles-zeros}

\begin{chapterlead}
The scattering matrix of an optical system encodes every linear response
property---reflection, transmission, absorption, and amplification---as a
function of frequency. Its singularity structure in the complex frequency
plane tells us when the system lases (a pole reaches the real axis), when
it perfectly absorbs (a zero sits on the real axis), and when it does
both simultaneously (the CPA-laser point). This pole--zero description
is used for non-Hermitian photonic systems through concrete examples.
It also establishes the connection between time reversal,
$\PT$ symmetry, and the CPA-laser effect~\cite{novitsky2019cpa,alizadeh2025lasing}, and clarifies the critical
distinctions between exceptional points, CPA-laser thresholds, and
maximum contrast conditions.
\end{chapterlead}


%% ================================================================
\section{Scattering-matrix singularities revisited}
\label{sec:singularities-revisited}

Chapter~\ref{ch:boundaries} introduced the poles and zeros of
$\Smat(\omega)$ as features of the complex frequency plane. We now
develop this picture more systematically.

\subsection{Poles: threshold lasing}
\label{ssec:lasing-poles}

A pole of $\Smat(\omega)$ at $\omtil = \omega_0 - i\gamma$ corresponds to
a quasinormal mode (Chapter~\ref{ch:qnm}). For a passive system,
$\gamma > 0$ and the pole lies below the real axis. If gain is introduced,
$\gamma$ decreases; at threshold, $\gamma \to 0$ and the pole reaches the
real axis:
\begin{equation}
 \Im\,\omtil_n = 0
 \qquad\Longleftrightarrow\qquad
 \text{lasing threshold}.
 \label{eq:lasing-threshold}
\end{equation}
At this point, a self-sustained oscillation exists at real frequency
$\omega_0$ with no external driving---the system is a laser.

The physical picture is transparent in the pole--zero language: gain lifts
the QNM poles toward the real axis. A passive cavity has all poles in the
lower half-plane; adding gain is equivalent to reducing the imaginary parts
of the poles. The first pole to reach the real axis sets the lasing
threshold and frequency.

\subsection{Zeros: coherent perfect absorption}
\label{ssec:cpa-zeros}

A zero of $\det\Smat(\omega)$ at real $\omega$ means there exists a
specific combination of incoming waves such that the outgoing field
vanishes identically. All incident energy is absorbed. This is
\emph{coherent perfect absorption} (CPA)~\cite{chong2010symmetry}.

CPA requires precise control of the relative phase and amplitude of the
incoming beams. For a two-port system, the CPA condition is
\begin{equation}
 \det\Smat(\omega_{\mathrm{CPA}}) = 0
 \quad\Longleftrightarrow\quad
 \text{an eigenvalue of } \Smat \text{ vanishes}.
 \label{eq:cpa-condition}
\end{equation}

\subsection{The meromorphic structure}
\label{ssec:meromorphic}

The scattering matrix $\Smat(\omega)$ of a finite photonic structure is a
meromorphic function of $\omega$: it is analytic everywhere in the complex
plane except at isolated poles. Any meromorphic matrix can be written in
terms of its poles and zeros as a product expansion. For a scalar
reflection or transmission coefficient $s(\omega)$, this takes the form
\begin{equation}
 s(\omega) = A\,\prod_j \frac{\omega - z_j}{\omega - p_j},
 \label{eq:pole-zero-product}
\end{equation}
where $\{p_j\}$ are the poles, $\{z_j\}$ are the zeros, and $A$ is a
constant prefactor. This representation makes the spectral response a
direct map of the singularity positions:
\begin{itemize}
 \item Near a pole $p_j$, the response $|s(\omega)|^2$ diverges---a
 sharp amplification peak.
 \item Near a zero $z_j$, the response vanishes---a sharp absorption dip.
 \item Between a nearby pole--zero pair, the response exhibits a Fano
 line shape: asymmetric with a rapid transition from absorption to
 amplification.
\end{itemize}


%% ================================================================
\section{Example: Fabry--P\'erot cavity}
\label{sec:fabry-perot-poles}

The Fabry--P\'erot cavity is the simplest photonic structure that
illustrates the full pole--zero framework.

\subsection{Setup and transfer matrix}
\label{ssec:fp-setup}

Consider a dielectric slab of thickness $L$, refractive index
$n = n' + in''$ (with $n'' > 0$ for loss and $n'' < 0$ for gain),
embedded in vacuum. The transfer matrix for the slab, including the two
air--dielectric interfaces, is (Appendix~\ref{app:transfer})
\begin{equation}
 \Tmat = \frac{1}{4n}\begin{pmatrix}
 (n+1)^2 e^{-i\phi} - (n-1)^2 e^{i\phi} & (n^2-1)(e^{i\phi}-e^{-i\phi}) \\
 (n^2-1)(e^{-i\phi}-e^{i\phi}) & (n+1)^2 e^{i\phi} - (n-1)^2 e^{-i\phi}
 \end{pmatrix},
 \label{eq:fp-transfer}
\end{equation}
where $\phi = n\omega L/c$ is the single-pass phase. The transmission
coefficient is $t = 1/T_{22}$.

\subsection{Pole positions}
\label{ssec:fp-poles}

The poles of $t(\omega)$ are the zeros of $T_{22}$, which satisfy the
round-trip resonance condition
\begin{equation}
 r_{12}^2\,e^{-2in\omega L/c} = 1,
 \label{eq:fp-roundtrip}
\end{equation}
where $r_{12} = (n-1)/(n+1)$ is the Fresnel reflection coefficient at
each interface. For weak material loss or gain, write
$r_{12} \approx |r_{12}|\,e^{i\varphi_r}$ using the real background
index $n'$. The pole frequencies are then approximately
\begin{equation}
 \omtil_m \approx \frac{c}{n' L}\bigl(m\pi - \varphi_r\bigr)
 - i\,\frac{c}{n' L}\,\ln\frac{1}{|r_{12}|},
 \qquad m = 0, \pm 1, \pm 2, \ldots
 \label{eq:fp-poles}
\end{equation}
For the lossless reference cavity, the imaginary part is negative: all
poles lie below the real axis. The mirror-limited decay rate is
$\gamma_m = (c/n'L)\ln(1/|r_{12}|)$; material absorption adds to this
decay, while material gain subtracts from it.

\paragraph{Counting poles and zeros.}
For the Fabry--P\'erot slab, the poles and zeros of the scalar
transmission $t(\omega)$ are approximately equally spaced in the complex
plane with a free spectral range $\Delta\omega = \pi c/(n'L)$ when the
dispersion and absorption are weak. In the lossless reference cavity,
zeros are the time reverses of the poles and appear at the conjugate
positions. Material loss or gain moves the zeros as well as the poles;
the zero motion is what makes coherent perfect absorption possible.

For a slab with $n' = 3.5$ and $L = 300\;$nm at telecom wavelengths
($\lambda \sim 1550\;$nm), the free spectral range is
$\Delta\omega \approx 4.49 \times 10^{14}\;$rad/s, and the Fresnel
reflectivity $|r_{12}|^2 = [(3.5-1)/(3.5+1)]^2 \approx 0.309$ gives
a quality factor $Q = \omega_0/(2\gamma_0) \approx 2.7$
for the lowest-order mode. The low $Q$ makes the pole--zero structure
well separated and easy to resolve, which is why this cavity is ideal
for pedagogical demonstration of the scattering singularity framework.

\subsection{Adding gain: pole motion toward lasing}
\label{ssec:fp-gain}

If the slab material has gain ($n'' < 0$), the imaginary part of the
propagation phase $\phi$ acquires a positive contribution that
counteracts the mirror loss. The pole frequency becomes
\begin{equation}
 \Im\,\omtil_m = -\frac{c}{n'L}\bigl[\ln(1/|r_{12}|) + n''\omega_m L/c\bigr].
 \label{eq:fp-pole-gain}
\end{equation}
When $|n''|\omega_m L/c = \ln(1/|r_{12}|)$, the imaginary part vanishes
and the pole reaches the real axis: this is the lasing threshold for mode
$m$. The threshold gain is
\begin{equation}
 g_{\mathrm{th}} = \frac{2\ln(1/|r_{12}|)}{L},
 \label{eq:fp-threshold}
\end{equation}
where $g_{\mathrm{th}} = 2|n''|\omega/c$ is the power gain coefficient.
This is the standard textbook result, but here it emerges naturally as a
statement about pole motion in the complex plane.

\subsection{Zero positions and CPA}
\label{ssec:fp-zeros}

The zeros of $\det\Smat$ for the Fabry--P\'erot slab satisfy the
time-reversed round-trip condition:
\begin{equation}
 r_{12}^{*2}\,e^{-2in^*\omega L/c} = 1.
 \label{eq:fp-zero-condition}
\end{equation}
For a purely lossy slab ($n'' > 0$), the zeros lie in the upper
half-plane. As loss increases, zeros descend toward the real axis.
CPA occurs when a zero reaches the real axis: the slab absorbs all
coherently incident light at the CPA frequency.

The CPA threshold loss is the mirror image of the lasing threshold gain:
\begin{equation}
 \alpha_{\mathrm{CPA}} = \frac{2\ln(1/|r_{12}|)}{L},
 \label{eq:fp-cpa-threshold}
\end{equation}
confirming the time-reversal duality~\cite{sweeney2019perfectly,wang2021coherent}: CPA is lasing run backward.

\paragraph{Explicit pole--zero map for the Fabry--P\'erot cavity.}
For the $n' = 3.5$, $L = 300\;$nm slab, the first three poles are at
$\omtil_m/\Delta\omega = m - i\,0.187$ for $m = 1, 2, 3$, and the
corresponding zeros at $\omtil_m^*/\Delta\omega = m + i\,0.187$. As
gain is added ($n'' \to -|n''|$), each pole moves vertically upward by
$\delta(\Im\omtil) = |n''|\omega_m/n'$. The lasing threshold for
mode $m$ is $|n''|_{\mathrm{th}} = \ln(1/|r_{12}|)\,c/(\omega_m L)
= 0.187\,\Delta\omega/\omega_m$. Simultaneously, each zero moves
downward by the same amount---the gain that brings a pole to the real
axis pushes its partner zero away. Only in a $\PT$-symmetric
configuration (gain and loss in separate halves of the cavity) do a
pole and zero converge on the real axis simultaneously, producing the
CPA-laser point.

%% ================================================================
\section{Time reversal and the CPA-laser duality}
\label{sec:time-reversal}

CPA and lasing are connected by time reversal. A laser emits coherent
radiation with no input; CPA absorbs coherent radiation with no output.
Formally, if $\Smat(\omega_L)$ has a pole (lasing), then
$\Smat^{-1}(\omega_L)$ has a zero---and $\Smat^{-1}$ is the scattering
matrix of the time-reversed system (with gain and loss
exchanged)~\cite{chong2010symmetry}.

\subsection{Formal statement}
\label{ssec:tr-formal}

Let $\Smat(\omega)$ be the scattering matrix of a linear system with
material parameters $\epsr(\rr)$. The time-reversed system has
$\varepsilon_{\mathrm{TR}}(\rr) = \epsr^*(\rr)$ (gain $\leftrightarrow$ loss),
and its scattering matrix is
\begin{equation}
 \Smat_{\mathrm{TR}}(\omega) = [\Smat(\omega^*)]^{*\,-1}.
 \label{eq:s-time-reversed}
\end{equation}
At real frequencies, $\Smat_{\mathrm{TR}}(\omega) = [\Smat(\omega)]^{-1\,*}$.
The poles of $\Smat$ become zeros of $\Smat_{\mathrm{TR}}$, and vice versa.

\subsection{CPA-laser point}
\label{ssec:cpa-laser}

For a system with both gain and loss, lasing and CPA can occur at the
\emph{same} frequency if and only if a pole and a zero of $\Smat$
coincide on the real axis. This simultaneous occurrence is the
\emph{CPA-laser} point.

A $\PT$-symmetric structure provides a natural setting for the CPA-laser
because the $\PT$ constraint relates the inverse scattering matrix to
the parity-transformed, time-reversed scattering matrix. For real
frequencies this is the two-port relation discussed in
Chapter~\ref{ch:pt-symmetric}; analytically continued into the complex
plane, it pairs a pole at $\omtil$ with a zero at $\omtil^*$.
When $\omtil$ reaches the real axis, the pole and zero meet
simultaneously---the CPA-laser condition is automatically satisfied.

%% ================================================================
\section{$\PT$-symmetric scattering}
\label{sec:pt-scattering}

For a $\PT$-symmetric structure, the scattering matrix satisfies the
fundamental relation~\eqref{eq:pt-s-relation}. At real frequencies, the
consequences are:
\begin{enumerate}
 \item $|\det\Smat| = 1$: the product of scattering eigenvalues has
 unit modulus.
 \item Poles and zeros of $\Smat$ come in conjugate pairs:
 if $\omtil$ is a pole, then $\omtil^*$ is a zero, and vice versa.
 \item In the $\PT$-symmetric phase, both scattering eigenvalues are
 unimodular: neither net amplification nor net absorption occurs.
 \item In the $\PT$-broken phase, one eigenvalue has $|s| > 1$
 (amplification) and the other has $|s| < 1$ (absorption).
\end{enumerate}

The CPA-laser point, where a pole and zero meet on the real axis,
corresponds to the special situation where $\PT$ symmetry enforces
$s_1 \to \infty$ and $s_2 \to 0$ simultaneously, keeping
$|s_1 s_2| = 1$.

\subsection{Scattering phase diagram}
\label{ssec:scattering-phase-diagram}

As the gain--loss parameter $\gamma$ increases from zero, the scattering
eigenvalues trace paths in the complex plane. Three qualitatively
different regimes emerge:
\begin{enumerate}
 \item \textbf{$\PT$-symmetric phase ($\gamma < \gamma_{\mathrm{EP}}$):}
 both $|s_1| = |s_2| = 1$. The system is transparent on average,
 though phase shifts differ from the passive case.
 \item \textbf{$\PT$-broken phase ($\gamma_{\mathrm{EP}} < \gamma < \gamma_{\mathrm{CPA}}$):}
 $|s_1| > 1$, $|s_2| < 1$. One channel amplifies, the other absorbs,
 but neither diverges.
 \item \textbf{CPA-laser threshold ($\gamma = \gamma_{\mathrm{CPA}}$):}
 $|s_1| \to \infty$, $|s_2| \to 0$. The system simultaneously lases
 and perfectly absorbs at the same frequency for different input states.
\end{enumerate}
This phase diagram highlights that the EP (a spectral degeneracy) and
the CPA-laser (a pole--zero coincidence) generally occur at
\emph{different} gain--loss values.

%% ================================================================
\section{Transfer-matrix analysis of a $\PT$ multilayer}
\label{sec:pt-multilayer}

\subsection{Setup}
\label{ssec:multilayer-setup}

Consider a multilayer stack of $N$ unit cells, each composed of a loss
layer (thickness $d_1$, refractive index $n_1 = n_0 + i\kappa_I$) and a
gain layer (thickness $d_2 = d_1$, refractive index
$n_2 = n_0 - i\kappa_I$). This structure satisfies the $\PT$ condition
$n(x) = n^*(-x)$ when the inversion center is placed at the interface
between the loss and gain layers.

\subsection{Output coefficients}
\label{ssec:output-coefficients}

The total transfer matrix $\Tmat^N$ determines the reflection and
transmission through~\eqref{eq:pt-transfer-to-scatter}. As a function of
the gain--loss parameter $\kappa_I$, three regimes
emerge~\cite{novitsky2019cpa}:
\begin{enumerate}
 \item \textbf{Below the EP} ($\kappa_I < \kappa_{\mathrm{EP}}$):
 the transmission and reflection show moderate asymmetry. The output
 intensity can exceed the input, but the contrast between amplification
 and absorption is limited.
 \item \textbf{At the EP} ($\kappa_I = \kappa_{\mathrm{EP}}$):
 the transfer-matrix eigenvalues coalesce. The output exhibits maximum
 contrast between amplification and absorption for specific input
 configurations.
 \item \textbf{Above the EP} ($\kappa_I > \kappa_{\mathrm{EP}}$):
 the system enters the $\PT$-broken phase. One scattering channel
 grows exponentially with $N$ (laser-like amplification), while the
 other is exponentially suppressed (CPA-like absorption).
\end{enumerate}

\subsection{CPA-laser versus exceptional point}
\label{ssec:cpa-vs-ep}

A common misconception is that the CPA-laser point coincides with the
exceptional point. In general, it does not~\cite{novitsky2019cpa}.
The EP is a property of the transfer matrix eigenvalues (a spectral
degeneracy); the CPA-laser is a property of the scattering matrix
singularities (a pole--zero coincidence on the real axis). The two
conditions are related but distinct:
\begin{itemize}
 \item The EP can occur at a gain--loss value where no real-frequency
 pole or zero exists.
 \item The CPA-laser point can occur in the broken phase, away from
 the EP.
 \item Maximum output contrast (the ratio of amplification to
 absorption for different input phases) need not coincide with
 either the EP or the CPA-laser point.
\end{itemize}

This distinction is physically important: the EP controls the spectral
topology and the eigenvector structure, while the CPA-laser controls the
input--output energy balance.

\begin{takeaway}[Three distinct critical points]
In a $\PT$-symmetric scattering system, the exceptional point (spectral
degeneracy), the CPA-laser point (pole--zero coincidence), and the point
of maximum amplification--absorption contrast are in general three
different parameter values. Conflating them leads to incorrect
predictions of device performance.
\end{takeaway}


%% ================================================================
\section{Spectral singularities}
\label{sec:spectral-singularities}

\subsection{Definition}
\label{ssec:spectral-sing-def}

A \emph{spectral singularity} is a point on the real frequency axis~\cite{longhi2010mml,chong2010symmetry}
where the reflection or transmission coefficient diverges---i.e., a pole
of $\Smat(\omega)$ at real $\omega$. Unlike the QNM poles discussed
above, which lie in the lower half-plane for passive systems, a spectral
singularity occurs on the real axis and is the scattering-theory
description of threshold lasing in a linear gain medium.

Mathematically, a spectral singularity at $\omega_s$ means the
resolvent of the Maxwell operator is unbounded at $\omega_s$, but
$\omega_s$ is not an eigenvalue in the standard sense: no
square-integrable eigenfunction exists. Instead, a spectral singularity
is a divergence of the scattering amplitude at a real frequency---an
``eigenvalue embedded in the continuum'' of the non-self-adjoint
operator.

\subsection{Physical interpretation}
\label{ssec:spectral-sing-physics}

For optical systems, spectral singularities have two complementary
physical manifestations that are connected by time reversal.

The first is \emph{lasing}. A pole of $\Smat(\omega)$ on the real
axis means that the system sustains a self-oscillating field at
frequency $\omega_s$ without any external drive. In the meromorphic
representation
\begin{equation}
 \Smat(\omega) = \Smat_{\mathrm{bg}}(\omega)\,
 \prod_n \frac{\omega - z_n}{\omega - p_n},
 \label{eq:smat-meromorphic}
\end{equation}
where $p_n$ are the QNM poles and $z_n$ the scattering zeros, lasing
occurs when one pole $p_m$ migrates to the real axis:
$\Im\,p_m \to 0^-$. At this point the denominator
$(\omega - p_m)$ vanishes for real $\omega = \Re\,p_m$, and the
scattering amplitudes diverge. The outgoing field grows without bound
in time---the signature of stimulated emission overcoming all losses.

The second manifestation is \emph{coherent perfect absorption} (CPA),
also called anti-lasing. A zero of $\det\Smat(\omega)$ on the real
axis means that there exists a coherent input field configuration that
is perfectly absorbed: no light exits the structure. Mathematically,
this is a spectral singularity of the \emph{inverse} scattering matrix
$\Smat^{-1}$. By~\eqref{eq:smat-meromorphic}, CPA occurs when a zero
$z_m$ sits on the real axis: $\Im\,z_m = 0$.

The connection between lasing and CPA follows from the time-reversal
structure of the S-matrix. A pole of a given system maps to a zero of
the time-reversed system, in which gain and loss are interchanged. In a
lossless reciprocal reference problem this reduces to the familiar
conjugate pole--zero symmetry, but it is not a general consequence of
reciprocity alone. In specially constrained active systems, most notably
$\PT$-symmetric ones, a pole and a zero can be forced to reach the real
axis together. This is the \emph{CPA-laser}
point~\cite{chong2010symmetry,sweeney2019perfectly}, at which the system
can act as either a laser or a perfect absorber depending on the input
coherence.

\paragraph{Quantitative threshold from S-matrix poles.}
Combining the meromorphic representation~\eqref{eq:smat-meromorphic}
with the Fabry--P\'erot pole positions from
Section~\ref{ssec:fp-poles}, the lasing threshold for mode $m$ is
reached when the gain $g$ satisfies
\begin{equation}
 g = g_{\mathrm{th},m}
 = -\frac{2\,\Im\,p_m(g=0)}{c/n_g}
 = \frac{2\ln(1/|r_{12}|)}{L}
 + \alpha_{\mathrm{mat}},
 \label{eq:threshold-from-poles}
\end{equation}
where $n_g$ is the group index, $\alpha_{\mathrm{mat}}$ accounts for
material absorption, and $|r_{12}|$ is the mirror reflectivity. The
first term is the mirror loss and the second the material loss; the
threshold condition requires the gain to compensate both. This
derivation makes explicit the pole-motion picture: gain ``lifts'' the
pole at a rate $c\,g/(2n_g)$ until it reaches the real axis.

\subsection{Threshold condition as a spectral singularity}
\label{ssec:spectral-threshold}

The spectral-singularity viewpoint provides a precise mathematical
criterion for the lasing threshold that goes beyond the simple
$\Im\omtil = 0$ condition. For a one-dimensional potential
$V(x) = -\omega^2[\varepsilon(x) - 1]/c^2$ with complex $\varepsilon$,
the spectral singularity occurs when the Jost function $f^+(\omega)$
vanishes at a real $\omega$ while $f^-(\omega) \neq 0$. This is a
stronger condition than a pole of the reflection coefficient: it
requires that the outgoing-wave solution be bounded (not
square-integrable) while the incoming-wave solution diverges.

For the Fabry--P\'erot cavity, the Jost-function condition reduces to
the round-trip condition~\eqref{eq:fp-roundtrip} evaluated on the real
axis. But for more complex structures---photonic crystals, coupled
resonators, random media---the Jost-function framework provides a
rigorous definition of the lasing threshold that does not require
identifying individual QNM poles.

\subsection{Spectral singularities in $\PT$-symmetric potentials}
\label{ssec:spectral-sing-pt}

In a $\PT$-symmetric optical potential, spectral singularities have
special properties. For reciprocal one-dimensional scattering, the
generalized conservation relation can be written as
\begin{equation}
 \bigl||t(\omega)|^2-1\bigr|=\sqrt{R_L(\omega)R_R(\omega)},
 \qquad R_{L,R}=|r_{L,R}|^2 .
 \label{eq:pt-generalized-unitarity}
\end{equation}
A spectral singularity at $\omega_s$ corresponds to
 $|t(\omega_s)| \to \infty$, and the reflection product diverges
simultaneously---the system emits coherent radiation into the available
channels. Its time-reversed partner is the CPA condition.


%% ================================================================
\section{Beyond $\PT$: general pole--zero engineering}
\label{sec:pole-zero-engineering}

The pole--zero framework extends beyond $\PT$-symmetric systems. For
any linear photonic structure described by a scattering matrix, the
motion of poles and zeros in the complex frequency plane under parameter
changes follows simple rules that enable systematic device design.
Gain moves poles upward in the complex plane, toward the real axis and
eventually through it to lasing. Loss moves poles downward, increasing
decay rates and broadening resonances. Structural changes---geometry,
coupling strength, grating period---move poles horizontally, shifting
the resonance frequency while leaving the decay rate largely unchanged.
The zero structure is constrained by causality, time-reversal relations,
and passivity bounds on real-frequency scattering, but zeros are not
fixed by reciprocity alone. Indeed, in passive absorbers one deliberately
moves a scattering zero onto the real axis by controlling the spatial
profile and strength of absorption.
These rules provide an intuitive design language: achieving CPA at a
target frequency amounts to placing a zero on the real axis, while
suppressing a parasitic lasing mode amounts to pushing its pole deeper
into the lower half-plane.

\subsection{Design as singularity placement}
\label{ssec:design-singularity}

Designing a photonic device---a laser, a sensor, a filter, or an
absorber---can be viewed as engineering the positions of poles and zeros
in the complex frequency plane to achieve the desired spectral response.
This perspective unifies many seemingly different problems:
\begin{itemize}
 \item A narrow-band laser requires one pole close to the real axis with
 all other poles far away.
 \item A broadband absorber requires zeros spread along the real axis.
 \item A Fano resonance requires a nearby pole--zero pair with
 controlled spacing.
 \item A flat-top filter requires poles and zeros arranged in a
 Butterworth-like pattern.
\end{itemize}

\paragraph{Metasurface example: pole--zero placement for perfect
absorption.}
Consider a metasurface consisting of a periodic array of lossy
dielectric resonators on a metallic back-reflector. The back-reflector
eliminates transmission, so the scattering problem reduces to a
single-port reflection coefficient $r(\omega)$. Perfect absorption at
a target frequency $\omega_0$ requires $r(\omega_0) = 0$, i.e., a zero
of $r$ on the real axis. In the pole--zero language, this is achieved
by tuning the resonator geometry so that a pole at
$\omtil = \omega_0 - i\gamma$ is paired with a zero at
$z = \omega_0 + i\gamma'$, and then adjusting the material loss so
that $\gamma' \to 0$. The critical-coupling condition
$\gamma_{\mathrm{rad}} = \gamma_{\mathrm{abs}}$ (radiative decay
equals absorptive decay) is precisely the statement that the zero
reaches the real axis.

With the global $e^{-i\omega t}$ convention used in this book, a
one-port temporal coupled-mode model may be written in the dimensionless
form
\begin{equation}
 r(\Delta)=
 \frac{-i\Delta+\gamma_{\mathrm{abs}}-\gamma_{\mathrm{rad}}}
 {-i\Delta+\gamma_{\mathrm{abs}}+\gamma_{\mathrm{rad}}},
 \qquad \Delta=\omega-\omega_0.
 \label{eq:one-port-critical-coupling}
\end{equation}
The pole and zero are therefore
\begin{equation}
 \omega_p=\omega_0-i(\gamma_{\mathrm{abs}}+\gamma_{\mathrm{rad}}),
 \qquad
 \omega_z=\omega_0+i(\gamma_{\mathrm{rad}}-\gamma_{\mathrm{abs}}).
 \label{eq:one-port-pole-zero}
\end{equation}
The passive zero starts in the upper half-plane when absorption is weak
and crosses the real axis at critical coupling.

\begin{figure}[t]
 \centering
 \includegraphics[width=0.96\linewidth]{fig_f10_critical_coupling_absorber.pdf}
 \caption{Critical coupling as zero placement in a one-port absorber.
 The time convention is $e^{-i\omega t}$ and the reflection amplitude is
 Eq.~\eqref{eq:one-port-critical-coupling}. (a)~Imaginary parts of
 the pole and zero in Eq.~\eqref{eq:one-port-pole-zero} as the ratio
 $\gamma_{\mathrm{abs}}/\gamma_{\mathrm{rad}}$ is varied. The pole
 remains in the lower half-plane, while the passive zero begins in the
 upper half-plane and reaches $\operatorname{Im}\omega_z=0$ at
 $\gamma_{\mathrm{abs}}=\gamma_{\mathrm{rad}}$. (b)~Real-frequency
 reflectance spectra for under-coupled
 ($\gamma_{\mathrm{abs}}/\gamma_{\mathrm{rad}}=0.5$), critically
 coupled, and over-coupled
 ($\gamma_{\mathrm{abs}}/\gamma_{\mathrm{rad}}=2$) cases. The two
 off-critical spectra have the same minimum reflectance,
 $|r(0)|^2=1/9$, while critical coupling gives $|r(0)|^2=0$.}
 \label{fig:ch10-critical-coupling-absorber}
\end{figure}

This design recipe generalises to multi-frequency perfect absorption:
each target frequency requires one zero on the real axis, which in turn
requires one resonator mode with matched radiative and absorptive
decay rates. The pole--zero framework provides a systematic way to
count the degrees of freedom and identify the minimum number of
resonator modes needed for a given absorption bandwidth.

There is an important physical distinction between critical coupling and
lasing threshold. Critical coupling is still a passive condition: the
pole remains in the lower half-plane because
$\operatorname{Im}\omega_p=-(\gamma_{\mathrm{abs}}+
\gamma_{\mathrm{rad}})<0$. The device absorbs perfectly at one real
frequency not because the stored field grows without bound, but because
the directly reflected wave and the field re-radiated by the resonator
cancel exactly in the outgoing channel. In Maxwell language, the
Poynting flux entering the resonator through the external channel equals
the Ohmic or material-loss power dissipated inside the resonator. This
is why the condition is often called \emph{impedance matching}: the
external radiation resistance and internal absorption resistance are
matched.

This clarification prevents a common misconception. Moving a zero to
the real axis is not the same as moving a pole to the real axis. The
first produces a finite, stable absorbing solution for a particular
coherent input; the second produces a self-sustained outgoing solution
with no input. The two events can coincide only in special active
systems such as the $\PT$-symmetric CPA-laser discussed above. In a
one-port passive metasurface, only the zero reaches the real axis, so
perfect absorption is compatible with stability, causality, and finite
stored energy.

\subsection{Passivity constraints}
\label{ssec:passivity-poles}

For a passive system (no gain), the poles of $\Smat(\omega)$ must lie in
the lower half-plane. This is a direct consequence of causality and
stability: the Green function of a passive system decays in the future,
not the past. Scattering zeros obey different constraints. In a lossless
unitary system they mirror the poles across the real axis; in a lossy
passive system they can move onto the real axis, producing CPA, or away
from the simple conjugate locations. Violating pole stability requires
gain, which is the fundamental reason why lasers and CPA-lasers are
active devices.

\subsection{Connection to QNM perturbation theory}
\label{ssec:pole-zero-qnm}

The biorthogonal perturbation theory of Chapter~\ref{ch:qnm} provides
the sensitivity of pole positions to material or geometric changes:
\begin{equation}
 \delta\omtil_n = -\frac{\omtil_n}{2}
 \frac{\int \delta\epsr\,\tilde{\EE}_n\cdot\tilde{\EE}_n\;\dd^3r}
 {\mathcal{N}_n}.
 \label{eq:pole-shift-perturbation}
\end{equation}
Here $\mathcal{N}_n$ is the unconjugated QNM normalization, including the
appropriate dispersive and surface or PML contribution. This formula
directly connects the pole--zero engineering picture to the underlying
Maxwell fields: placing gain where the relevant QNM has large overlap
moves the pole most efficiently toward the real axis.


%% ================================================================
\section{Connections}
\label{sec:ch10-forward}

This closes the scattering part of the book. Exceptional points,
$\PT$-symmetric optics, and pole--zero motion describe complementary
aspects of non-Hermitian photonic systems driven by gain, loss, and
coupling.

Part~IV now turns to the topological structure of non-Hermitian photonic
bands. Chapter~\ref{ch:berry-maxwell} builds Berry curvature directly
from Maxwell fields; Chapter~\ref{ch:complex-bands} introduces complex
bands and non-Hermitian gaps; and the subsequent chapters develop the
photonic skin effect, point-gap topology, and the reconstruction of
bulk-boundary correspondence.
